\documentclass[10pt,a4paper]{amsart}
\usepackage[margin=19mm]{geometry}
\usepackage{amsmath,amssymb,mathtools}
\usepackage[scr=rsfs,cal=euler]{mathalfa}
\usepackage{xfrac}
\usepackage[unicode,hidelinks,bookmarks=false]{hyperref}
\newcommand{\ie}{i.e.\ }
\newcommand{\cf}{cf.\ }
\newcommand{\resp}{respectively\ }
\newcommand{\Subsection}{\subsection}
\providecommand{\nobreakdash}{\nobreak\hbox{-}}
\setlength{\emergencystretch}{2em}
\multlinegap0pt
\usepackage{bm}
\usepackage{bbm}
\usepackage{dsfont}
\usepackage{xspace}
\usepackage{cancel}
\usepackage{mathtools}
\usepackage{amsthm}
\makeatletter
\@ifundefined{theorem}{\newtheorem{theorem}{Theorem}}{}
\makeatother
\title[Asymptotic-preserving conservative semi-Lagrangian discontin]{Asymptotic-preserving conservative semi-Lagrangian discontinuous Galerkin schemes for the Vlasov-Poisson system in the quasi-neutral limit}
\author{Xiaofeng Cai, Linghui Kong, Dmitri Kuzmin, Li Shan}
\date{Source: arXiv:2510.15854v1 (CC BY 4.0)}
\begin{document}
\maketitle

\noindent\emph{Source and licence.} Asymptotic-preserving conservative semi-Lagrangian discontinuous Galerkin schemes for the Vlasov-Poisson system in the quasi-neutral limit. Version arXiv:2510.15854v1 (CC BY 4.0), distributed under the Creative Commons Attribution 4.0 International licence, as recorded in the arXiv licence metadata for this exact version and shown on the abstract page https://arxiv.org/abs/2510.15854v1. Full rights notice: licenses/arxiv-2510.15854-CC-BY-4.0.txt.\par\smallskip

\noindent\textit{Editorial scope.} Counted unit: the Theorem whose source label is \texttt{None} in the article (source0.tex lines 946--960, source label \texttt{None}), delivered together with the article's own complete original proof of it (source0.tex lines 961--1074, one \texttt{proof} environment of 114 lines). No mathematics is written, repaired or re-derived: statement and proof are the article's own text line for line. Required context retained uncounted: HhE (lines 683--685). Every definition, display and label of those lines is retained unchanged. Omitted from this package and not redistributed: 1--682, 686--945, 1075--1945. Nothing inside a retained excerpt is omitted or altered: every retained body is delivered exactly as the article's lines are written, so no omission receipt of any kind accompanies this package. Citations: the retained excerpts carry no citation command at all, so no bibliography entry of the article is transcribed and no reference is redistributed. Reference adaptation: none was needed and none was made; every reference inside the delivered unit resolves inside it, so no unresolved reference or citation is produced. Printed slips kept literal: no printed word, symbol, hypothesis or number of the counted statement or of its proof was altered. Layout and class: the article's own class is \texttt{article} with packages including opticameet3, multirow, multicol, amsfonts, booktabs, pifont, amsthm, amsmath, cases, amssymb, graphicx, subcaption, url, caption; the wrapper is the lane's pinned \texttt{amsart} editorial wrapper at 10pt on A4, which restates the theorem-like environments the retained text needs and adds the article's own macro definitions that the retained text needs (none), with extra wrapper packages (bm, bbm, dsfont, xspace, cancel, mathtools). The delivered numbering is the unit's own. Assets: no retained span contains a figure environment, a graphics-inclusion command, a PDF-page inclusion, a table or a TikZ picture; no image file is copied into this package, the package declares no asset dependency and no image file is redistributed with it. Licence: version arXiv:2510.15854v1 of this article is distributed under the Creative Commons Attribution 4.0 International licence, as recorded in the arXiv licence metadata for this exact version and shown on the abstract page https://arxiv.org/abs/2510.15854v1. A full rights notice is delivered at licenses/arxiv-2510.15854-CC-BY-4.0.txt. Characters: every retained excerpt is pure ASCII, so the delivered engine, XeLaTeX with its default Latin Modern text font, reports no missing character.
\begin{align}\label{HhE}
    \int_{I_{v_i}} f_h^{n+1}(x_{j,q},v) \Psi(v)\mathrm d v=\int_{I_{v_i}^*} f_h^{*}(x_{j,q},v) \psi^n(v)\mathrm d v
\end{align}

\begin{theorem}[Linear stability of the operator $\mathcal{S}_{\mathbf{E}}^h$]
 The update $f_h^{n+1}=\mathcal{S}_{\mathbf{E}}^h(\Delta t)f_h^*$ using
 periodic boundary conditions and the linearization
 of the system $H_{\mathbf E}$ about the steady state
\begin{align*}
    \left\{
        \begin{aligned}
            &\rho=1,\\
            &\rho u=0,\\
            &\partial_x\phi=0
        \end{aligned}
    \right.
\end{align*}
is unconditionally $L^2$ stable in the sense of von Neumann stability analysis.
\end{theorem}

\begin{proof}
Using $\Psi(v)=1$ in \eqref{HhE}, summing the resulting equations
\begin{align}
    \int_{I_{v_i}}f_h^{n+1}(x_{j,q},v)\mathrm dv=\int_{I_{v_i}^*}f_h^*(x_{j,q},v)\mathrm dv
\end{align}
 over $i=1,\ldots,N_v$ and using the assumption of periodicity, we obtain
\begin{align}\label{0m}
    \int_{\Omega_v}f_h^{n+1}(x_{j,q},v)\mathrm dv=\int_{\Omega_v}f_h^*(x_{j,q},v)\mathrm dv.
\end{align}
This integral identity can be written as
\begin{align}\label{rhoequality}
    \rho^{n+1}(x_{j,q})=\rho^*(x_{j,q}).
\end{align}
Regardless of the mesh $\mathcal{E}_x$ for the discretization of $\Omega_x$, equation \eqref{rhoequality} is valid at all Gaussian quadrature points $x_{j,q}$. Summation over the indices of these points reveals that $\rho^{n+1}(x)=\rho^*(x)$ a.e. in
$\Omega_x$.

Next, we use $\Psi(v)=v$ in \eqref{HhE} to show that
\begin{align}
    \int_{I_{v_i}}f_h^{n+1}v\mathrm dv&=\int_{I_{v_i}^*}f_h^*(v+\Delta t\partial_x\phi^{n+1})\mathrm dv\notag\\
    &=\int_{I_{v_i}^*}f_h^*v\mathrm dv+\Delta t\partial_x\phi^{n+1}\int_{I_{v_i}^*}f_h^*
    \mathrm dv.
\end{align}
Summation 
over $i=1,\ldots,N_v$ under the assumption of periodicity gives
\begin{align}\label{1m}
    \int_{\Omega_v}f_h^{n+1}v\mathrm dv=\int_{\Omega_v}f_h^*v\mathrm dv
    +\Delta t\partial_x\phi^{n+1}\int_{\Omega_v}f_h^*\mathrm dv.
\end{align}
This proves that
$\frac{(\rho u)^{n+1}-(\rho u)^*}{\Delta t}-\rho^*\partial_x\phi^{n+1}=0$
a.e. in $\Omega_x$.

Combining the above auxiliary results, we consider the nonlinear semi-discrete problem
\begin{align}\label{nonlinsysHE}
    \left\{
        \begin{aligned}
            &\rho^{n+1}-\rho^*=0,\\
            &\frac{(\rho u)^{n+1}-(\rho u)^*}{\Delta t}-\rho^*\partial_x\phi^{n+1}=0,\\
            &\lambda^2\frac{\partial_{xx}\phi^{n+1}-2\partial_{xx}\phi^n+\partial_{xx}\phi^{n-1}}{\Delta t^2}+\partial_x(\rho^*\partial_x\phi^{n+1})=\partial_{xx}S^n,
        \end{aligned}
    \right.
\end{align}
where $S=\rho u^2+p$ and $p=\rho^{\gamma}$ with $\gamma>1$ is the pressure.
Linearizing the system \eqref{nonlinsysHE} about the steady state $\rho=1,~\mathcal{J}=\rho u=0,~\partial_x\phi=0$, we arrive at
\begin{align}
    \left\{
        \begin{aligned}
            &\rho^{n+1}-\rho^*=0,\\
            &\frac{\mathcal{J}^{n+1}-\mathcal{J}^*}{\Delta t}-\partial_x\phi^{n+1}=0,\\
            &\lambda^2\frac{\partial_{xx}\phi^{n+1}-2\partial_{xx}\phi^n+\partial_{xx}\phi^{n-1}}{\Delta t^2}+\partial_{xx}\phi^{n+1}=\partial_{xx}\rho^*.
        \end{aligned}
    \right.
\end{align}

Applying the spatial Fourier transform to $\rho,~\mathcal{J}$ and $\phi$, we denote the transformed variables by $\hat\rho,~\hat{\mathcal{J}}$ and $\hat\phi$, respectively. Introducing $\hat\psi^{n+1}=(\hat\phi^{n+1}-\hat\phi^n)/\Delta t$, we write
the linear system
\begin{align}
    \left\{
        \begin{aligned}
            &\hat\rho^{n+1}-\hat\rho^*=0,\\
            &\hat {\mathcal{J}}^{n+1}-\hat{\mathcal{J}}^*-ik\Delta t\hat\phi^{n+1}=0,\\
            &\hat\phi^{n+1}-\hat\phi^n-\Delta t\hat\psi^{n+1}=0,\\
            &\lambda^2\hat\psi^{n+1}-\lambda^2\hat\psi^n+\Delta t\hat\phi^{n+1}-\Delta t\hat\rho^*=0
        \end{aligned}
    \right.
\end{align}
in the matrix form
\begin{align}
B
    \left(
    \begin{array}{cccc}
    \hat\rho^{n+1}  \\
    \hat {\mathcal{J}}^{n+1}  \\
    \hat\phi^{n+1}   \\
    \hat\psi^{n+1}  \\
    \end{array}
\right)=A
    \left(
    \begin{array}{cccc}
    \hat\rho^*  \\
    \hat {\mathcal{J}}^*  \\
    \hat\phi^n  \\
    \hat\psi^n  \\
    \end{array}
\right),
\end{align}
where
\begin{align*}
B=\left(
\begin{array}{cccc}
    1 & 0 & 0 & 0 \\
    0 & 1 & -ik\Delta t & 0 \\
    0 & 0 & 1 & -\Delta t \\
    0 & 0 & \Delta t & \lambda^2 \\
\end{array}
\right),\quad A=\left(
\begin{array}{cccc}
    1 & 0 & 0 & 0 \\
    0 & 1 & 0 & 0 \\
    0 & 0 & 1 & 0 \\
    \Delta t & 0 & 0 & \lambda^2 \\
\end{array}
\right).
\end{align*}
To prove stability, we need to show that the modulus of each complex-valued
eigenvalue $\mu$ of $B^{-1}A$
is less than or equal to 1. We easily deduce that
$\mu_1=\mu_2=1$ and $\mu_3=\frac{\lambda^2+i\lambda\Delta t}{\lambda^2+\Delta t^2}$, $\mu_4=\frac{\lambda^2-i\lambda\Delta t}{\lambda^2+\Delta t^2}$. Since
\begin{align*}
    |\mu_3|=|\mu_4|=\sqrt{\frac{\lambda^4+\lambda^2\Delta t^2}{\lambda^4+2\lambda^2\Delta t^2+\Delta t^4}}<1,
    \end{align*}
 the linearized scheme satisfies von Neumann's $L^2$ stability
    criterion unconditionally.
\end{proof}

\end{document}
